\subsubsection{Sextic order, \texorpdfstring{$O(\gs^4)$}{O(gs**4)}}
\label{p:sec:hs:sextic}

The sextic order is the first where the auxiliary field of the previous order induces a composite vertex to be decomposed.
The vertex to add is therefore $\vertex{6}$ minus the exchange term of \cref{p:sec:hs:quintic}.
In its simplest implementation, the $3$-string field $\Sigma_3$ introduced at this order has a kinetic operator which is not a vertex of closed string field theory.


\paragraph{Decomposition of the reduced vertex}
\label{p:sec:hs:sextic:vertex}

The target at this order is:
\begin{equation}
  S_{\le 6}
    = S_{\le 5} + \frac{\gs^4}{6!} \vertex{6}(\Psi^6).
\end{equation}
The cubic action $\caction[\le 5]$ of the quintic order induces a new term beyond $S_{\le 5}$: by~\eqref{p:eq:hs5:eliminate}, eliminating $\Sigma_2$ from it produces the six-point term $\tfrac{\gs^4}{600} \vertex{5 \cdot 5}(\Psi^6)$, the exchange of $\Sigma_2$ between two quintic vertices.
The sextic order must thus reproduce on-shell the difference:
\begin{equation}
  \label{p:eq:hs6:remainder}
  S_{\le 6} - \caction[\le 5] \big|_{\Sigma_2 = \Sigma_2^{(3)*}}
    = \frac{\gs^4}{6!} \Big[ \vertex{6}(\Psi^6) - \frac{6}{5} \, \vertex{5 \cdot 5}(\Psi^6) \Big]
      + O(\gs^5).
\end{equation}
Two things have to happen at once: the sextic vertex must be produced, and the induced exchange term must be cancelled.
The second term is not a vertex of closed string field theory but a diagram built from $\vertex{4}$ and $\vertex{5}$.

We define the \emph{reduced vertex} $\dvertex{6}(A_1, \ldots, A_6)$ with general arguments as the fundamental vertex minus the composite vertex ~\eqref{p:eq:hs5:exchange-vertex}:
\begin{equation}
  \label{p:eq:hs6:dressed-vertex}
  \dvertex{6}(A_1, \ldots, A_6)
    := \vertex{6}(A_1, \ldots, A_6) - \frac{6}{5} \, \vertex{5 \cdot 5}(A_1, \ldots, A_6).
\end{equation}
This is represented graphically in \cref{p:fig:hs6:dressed-vertex}.
It is a six-linear form, graded-symmetric within each triple of arguments and under the exchange of the two triples, which is the symmetry of its exchange part.
It is not symmetric under the permutations mixing the two triples.
Read as a form in three $2$-string arguments, $\dvertex{6}(B_1, B_2, B_3)$ with $B_i \in \Hn{2}$, it is symmetrized in them, as $\vertex{5 \cdot 5}$ is in \cref{p:sec:hs:quintic}.
On the diagonal, where every argument is built from $\Psi$ alone, all readings agree: $\dvertex{6}(\Psi^3, \Psi^3) = \dvertex{6}(\Psi^2, \Psi^2, \Psi^2) = \dvertex{6}(\Psi^6)$.
By~\eqref{p:eq:hs6:remainder}, the term to add on-shell is thus the sextic term with $\dvertex{6}(\Psi^6)$ in place of $\vertex{6}(\Psi^6)$.
\begin{figure}[htp]
  \centering
  \begin{tikzpicture}
    \begin{feynman}
      \vertex (w) at (0, 0) [dot, label=above:{$\mathcal{W}_6$}];
      \vertex (w1) at (-0.9, 0.8) {};
      \vertex (w2) at (-1.1, 0) {};
      \vertex (w3) at (-0.9, -0.8) {};
      \vertex (w4) at (0.9, 0.8) {};
      \vertex (w5) at (1.1, 0) {};
      \vertex (w6) at (0.9, -0.8) {};
      \vertex (eq) at (1.9, 0) {$=$};
      \vertex (v) at (3.3, 0) [dot, label=above:{$\mathcal{V}_6$}];
      \vertex (v1) at (2.4, 0.8) {};
      \vertex (v2) at (2.2, 0) {};
      \vertex (v3) at (2.4, -0.8) {};
      \vertex (v4) at (4.2, 0.8) {};
      \vertex (v5) at (4.4, 0) {};
      \vertex (v6) at (4.2, -0.8) {};
      \vertex (minus) at (5.3, 0) {$- \; \dfrac{6}{5}$};
      \vertex (a) at (7.0, 0) [dot];
      \vertex (b) at (9.0, 0) [dot];
      \vertex (a1) at (6.2, 0.8) {};
      \vertex (a2) at (6.0, 0) {};
      \vertex (a3) at (6.2, -0.8) {};
      \vertex (b1) at (9.8, 0.8) {};
      \vertex (b2) at (10.0, 0) {};
      \vertex (b3) at (9.8, -0.8) {};
      \diagram* {
      (w) -- (w1), (w) -- (w2), (w) -- (w3), (w) -- (w4), (w) -- (w5), (w) -- (w6),
      (v) -- (v1), (v) -- (v2), (v) -- (v3), (v) -- (v4), (v) -- (v5), (v) -- (v6),
      (a) -- (a1), (a) -- (a2), (a) -- (a3), (b) -- (b1), (b) -- (b2), (b) -- (b3),
      (a) -- [plain, double, edge label={$\Sigma_2$}] (b),
      };
    \end{feynman}
  \end{tikzpicture}
  \caption{The reduced sextic vertex~\eqref{p:eq:hs6:dressed-vertex}: the sextic vertex minus the composite vertex, two quintic vertices joined by a $\Sigma_2$ line.}
  \label{p:fig:hs6:dressed-vertex}
\end{figure}


We have several options to decompose this vertex, since the matching constrains the sextic term only through its value on the diagonal, which is a sextic polynomial in $\Psi$.
Any trilinear form on $\Hn{2}$ with the value $\dvertex{6}(\Psi^6)$ on $\Sigma_2 = \Psi^2$ reproduces the target at this order.
Two forms differing by one that vanishes on the diagonal give the same eliminated action: they are equivalent at this truncation.
The construction we keep is the one that stays exact at all orders.
We see that we cannot immediately apply the trick of the quartic order with a $3$-string field $\Sigma_3 \in \Hn{3}$, because we would get a quartic term $\mc V_6(\Sigma_3, \Psi^3)$.
Another solution could be to write $\mc V_6(\Sigma_2^3)$, but this would not generalize at higher orders since $\mc V_8(\Sigma_2^4)$ is quartic.\footnotemark{}
\footnotetext{
  Such a term appeared in~\cite{Erbin:2024:OpenStringStub}.
  This term together with the terms linear and quadratic in the auxiliary field were sufficient to reproduce the complete action, with a single auxiliary field.
  This is specific to the structure of the interactions induced by the stubs.
}
The way out is to take an intermediary path and replace $\Psi^2$ by $\Sigma_2$ in the tensor power, since the solution $\Sigma_2 = \Psi^2 + O(\gs)$ guarantees that only $\dvertex{6}(\Psi^6)$ is produced at this order.
Moreover, this form is suggested by the form of the composite vertex, whose arguments are identical, $\vertex{5 \cdot 5}(\Sigma_2^2, \Psi^2)$.


We write the sextic term as $\dvertex{6}(\Psi \otimes \Sigma_2, \Psi \otimes \Sigma_2)$, with the $3$-string state $\Psi \otimes \Sigma_2$ in each triple of slots, and add the square of the fluctuation of $\Sigma_3$ from it:
\begin{equation}
  \label{p:eq:hs6:DeltaS6}
  \Delta S_6
    := - \frac{\gs^4}{6!} \, \dvertex{6}(\Sigma_3 - \Psi \otimes \Sigma_2, \Sigma_3 - \Psi \otimes \Sigma_2),
\end{equation}
where
\begin{equation}
  \Sigma_3 \in \Hn{3},
  \qquad
  \Ngh(\Sigma_3)
    = 6.
\end{equation}
By construction, $\Delta S_6$ vanishes when $\Sigma_3 = \Psi \otimes \Sigma_2$.
Expanding the square, we get the sextic term of the cubic action:
\begin{equation}
  \label{p:eq:hs6:S6tilde}
  \begin{aligned}
    \caction[6]
      &
      := \frac{\gs^4}{6!} \, \dvertex{6}(\Psi \otimes \Sigma_2, \Psi \otimes \Sigma_2) + \Delta S_6
      \\ &
      = - \frac{\gs^4}{6!} \, \dvertex{6}(\Sigma_3^2)
        + \frac{2 \gs^4}{6!} \, \dvertex{6}(\Psi \otimes \Sigma_2, \Sigma_3),
  \end{aligned}
\end{equation}
a quadratic term for $\Sigma_3$ and a cubic coupling of $\Psi$, $\Sigma_2$ and $\Sigma_3$.
As a consequence, the action through this order,
\begin{equation}
  \label{p:eq:hs6:action-tilde}
  \caction[\le 6]
    := \caction[\le 5] + \caction[6],
\end{equation}
is cubic in $(\Psi, \Sigma_2, \Sigma_3)$.
In terms of the elementary vertices of closed string field theory, we get:
\begin{equation}
  \label{p:eq:hs6:S6tilde-explicit}
  \begin{aligned}
    \caction[6]
      =
        &
        - \frac{\gs^4}{6!} \, \vertex{6}(\Sigma_3^2)
        + \frac{\gs^4}{600} \, \vertex{5}\Big( \Sigma_3, \sprod{2}{2}^{-1} \circ \sprod{3}{2}(\Sigma_3) \Big)
      \\ &
      + \frac{2 \gs^4}{6!} \, \vertex{6}(\Psi, \Sigma_2, \Sigma_3)
        - \frac{2 \gs^4}{600} \, \vertex{5}\Big( \Sigma_3, \sprod{2}{2}^{-1} \circ \sprod{3}{2}(\Psi, \Sigma_2) \Big),
  \end{aligned}
\end{equation}
The second and fourth terms, the kinetic term of $\Sigma_3$ and its coupling to $(\Psi, \Sigma_2)$, are both non-local in the sense of effective field theory, since they contain the inverse of a vertex, $\sprod{2}{2}^{-1}$.
This is the price of cancelling an exchange diagram with a contact term, and a typical phenomenon in effective field theory.
For generic vertices, in which the exchange term is not a linear combination of the elementary vertices, the elementary vertices alone cannot cancel the exchange term, which contains $\sprod{2}{2}^{-1}$: the subtraction must itself be an exchange term, produced either by the reduced vertex or by a second $2$-string field.
An alternative is to enlarge the field content with a second $2$-string field.
  We show it at the sextic order in \cref{p:app:wrong-sign} and will not develop this approach further.

Since the sextic vertex is now read with three of its slots as outputs, we define the product $\sprod{3}{3} \colon \Hn{3} \to \Hn{3}$ by:
\begin{equation}
  \label{p:eq:hs6:p33}
  \vertex{6}(A_1, \ldots, A_6)
    =: \pair[3]{A_1 \otimes A_2 \otimes A_3}{\sprod{3}{3}(A_4, A_5, A_6)},
\end{equation}
with the $3$-string pairing on the left, as $\sprod{2}{2}$ and $\sprod{3}{2}$ were defined from $\vertex{4}$ and $\vertex{5}$.
It has ghost number three and it is symmetric with respect to the $3$-string pairing since $\vertex{6}$ is graded-symmetric.


\paragraph{Elimination of the auxiliary fields}
\label{p:sec:hs:sextic:eom}

The $3$-string field enters $\caction[\le 6]$ only through $\dvertex{6}$.
We write this form through the $3$-string pairing:
\begin{equation}
  \label{p:eq:hs6:kinetic}
  \begin{aligned}
    \dvertex{6}(A, B)
      &
      =: \pair[3]{A}{\wprod{3}{3}(B)},
    \\
    \wprod{3}{3}
      &
      := \sprod{3}{3} - \frac{6}{5} \, \sprod{2}{3} \circ \sprod{2}{2}^{-1} \circ \sprod{3}{2}
    \qquad
    \wprod{3}{3} \colon \Hn{3} \to \Hn{3},
  \end{aligned}
\end{equation}
for $A, B \in \Hn{3}$, where $\sprod{2}{3} \colon \Hn{2} \to \Hn{3}$ is the transpose of $\sprod{3}{2}$,~\eqref{p:eq:conv:transpose}.
This is the kinetic operator of $\Sigma_3$.
It is symmetric with respect to the $3$-string pairing, as $\sprod{3}{3}$ and $\sprod{2}{2}^{-1}$ are, and it does not depend on the fields.
It receives no correction at higher orders.

Varying~\eqref{p:eq:hs6:S6tilde} with respect to the three fields, we get:
\begin{equation}
  \label{p:eq:hs6:delta-S}
  \begin{aligned}
    \delta \caction[6]
      &
      = \frac{2 \gs^4}{6!} \, \dvertex{6}(\delta \Psi \otimes \Sigma_2, \Sigma_3)
        + \frac{2 \gs^4}{6!} \, \dvertex{6}(\Psi \otimes \delta \Sigma_2, \Sigma_3)
      \\ &
      \quad
      - \frac{2 \gs^4}{6!} \, \pair[3]{\delta \Sigma_3}{\wprod{3}{3}(\Sigma_3)}
        + \frac{2 \gs^4}{6!} \, \pair[3]{\delta \Sigma_3}{\wprod{3}{3}(\Psi \otimes \Sigma_2)}.
  \end{aligned}
\end{equation}
The equations of motion of $\caction[\le 6]$ are:
\begin{subequations}
  \label{p:eq:hs6:eom}
  \begin{align}
    \label{p:eq:hs6:eom-psi}
    \mathcal E^{(4)}_\Psi
      &
      := \mathcal E^{(3)}_\Psi
        + \frac{2 \gs^4}{6!} \, \dvertex{6}(\, \cdot \, \otimes \Sigma_2, \Sigma_3)
      = 0,
    \\
    \label{p:eq:hs6:eom-sigma2}
    \mathcal E^{(4)}_{\Sigma_2}
      &
      := \mathcal E^{(3)}_{\Sigma_2}
        + \frac{\gs^2}{30} \, \dvertex{6}(\Psi \otimes \cdot \,, \Sigma_3)
      = 0,
    \\
    \label{p:eq:hs6:eom-sigma3}
    \mathcal E^{(4)}_{\Sigma_3}
      &
      := \wprod{3}{3}(\Psi \otimes \Sigma_2) - \wprod{3}{3}(\Sigma_3)
      = 0,
  \end{align}
\end{subequations}
where $\mathcal E^{(3)}_\Psi$ and $\mathcal E^{(3)}_{\Sigma_2}$ are the equations~\eqref{p:eq:hs5:eom} of the quintic order, and where $\dvertex{6}(\, \cdot \, \otimes \Sigma_2, \Sigma_3)$ denotes the state paired with the open slot.

The equation for $\Sigma_3$ is algebraic and linear, with the source $\wprod{3}{3}(\Psi \otimes \Sigma_2)$ and the field-independent kinetic operator $\wprod{3}{3}$.
A kernel direction $N_3 \in \ker \wprod{3}{3}$ appears neither in the quadratic term nor in the coupling of~\eqref{p:eq:hs6:S6tilde}, by the symmetry of $\wprod{3}{3}$.
It has zero action and couples to nothing at this order, as the kernel of $\sprod{2}{2}$ at the quartic order.
We therefore take the $3$-string field on the quotient:
\begin{equation}
  \label{p:eq:hs6:quotient}
  \Sigma_3
    \in \Hn{3} / \ker \wprod{3}{3},
\end{equation}
where $\wprod{3}{3}$ is injective by construction.
The equation of motion~\eqref{p:eq:hs6:eom-sigma3} then has the unique solution:
\begin{equation}
  \label{p:eq:hs6:saddle-sigma3}
  \Sigma_3^{(4)*}
    = \Psi \otimes \Sigma_2,
\end{equation}
for any $\Sigma_2$, and $\Psi \otimes \Sigma_2^{(4)*}$ on the joint solution.
The $3$-string field satisfies the reality condition $\Sigma_3^\star = - \Sigma_3$, the sign of three star-odd strands, and the kernel of $\wprod{3}{3}$ is stable under $\star$, so that the condition is meaningful on the quotient~\eqref{p:eq:hs6:quotient}.
The solution is real for real $\Psi$ and $\Sigma_2$ (\cref{p:app:reality}).
The $3$-string field is thus an auxiliary field: it is fixed by $(\Psi, \Sigma_2)$ through an algebraic equation, and its solution lies on the diagonal, $\Sigma_3 = \Psi^3 + O(\gs)$ once $\Sigma_2$ is on-shell.
Since $\Delta S_6$ is the square of $\Sigma_3 - \Psi \otimes \Sigma_2$, it vanishes on the solution.
Eliminating $\Sigma_3$ thus gives back the sextic term we started from, exactly in $\gs$ and for any $\Sigma_2$, as in~\eqref{p:eq:hs4:eliminate}:
\begin{equation}
  \label{p:eq:hs6:eliminate-sigma3}
  \caction[\le 6] \big|_{\Sigma_3 = \Sigma_3^{(4)*}}
    = \caction[\le 5] + \frac{\gs^4}{6!} \, \dvertex{6}(\Psi \otimes \Sigma_2, \Psi \otimes \Sigma_2).
\end{equation}

The elimination of $\Sigma_2$ then proceeds as at the quintic order, with one new ingredient.
The sextic term is quadratic in $\Sigma_2$ and of order $\gs^4$, so it adds a correction of order $\gs^2$ to the dressed operator of $\Sigma_2$ and none to its source, which remains $\sprod{2}{2}(\Psi^2)$.
The solution therefore becomes $\Sigma_2^{(4)*} = \Sigma_2^{(3)*} + O(\gs^2)$.
By stationarity of $\caction[\le 5]$ at $\Sigma_2^{(3)*}$, as in \cref{p:sec:hs:quintic}, a shift of order $\gs^2$ changes $\caction[\le 5]$ at order $\gs^2 \times \gs^4$, and the sextic term, of order $\gs^4$, at the same order.
At order $\gs^4$ only the value of the sextic term on $\Sigma_2 = \Psi^2$ enters, $\dvertex{6}(\Psi^6)$.
With~\eqref{p:eq:hs5:eliminate} and~\eqref{p:eq:hs6:remainder}, we get:
\begin{equation}
  \label{p:eq:hs6:eliminate}
  \begin{aligned}
    \caction[\le 6] \big|_{\Sigma_3 = \Sigma_3^{(4)*}, \, \Sigma_2 = \Sigma_2^{(4)*}}
      &
      = S_{\le 5}
        + \frac{\gs^4}{600} \, \vertex{5 \cdot 5}(\Psi^6)
        + \frac{\gs^4}{6!} \Big[ \vertex{6}(\Psi^6) - \frac{6}{5} \, \vertex{5 \cdot 5}(\Psi^6) \Big]
        + O(\gs^5)
      \\ &
      = S_{\le 6} + O(\gs^5).
  \end{aligned}
\end{equation}
The exchange term produced by $\Sigma_2$ at the quintic order and the exchange part of the reduced vertex cancel exactly, and the sextic vertex is reproduced.
The elimination of $\Sigma_3$ is exact, and the elimination of $\Sigma_2$ matches the target through order $\gs^4$.
Its remainder of order $\gs^5$ is the next induced interaction, built from the reduced vertex and from chains of quintic vertices on the $\Sigma_2$ line, which the septic order must absorb in its turn (\cref{p:sec:hs:septic}).
Off-shell, $\caction[\le 6]$ contains the square of the fluctuation $\Sigma_3 - \Psi \otimes \Sigma_2$ with the field-independent operator $\wprod{3}{3}$.
Once $\Sigma_3$ is eliminated, what is left is quadratic in $\Sigma_2$, with the dressed operator $\kin{2}{\Psi}{3}$ corrected at order $\gs^2$ by the sextic term.
The gauge transformations built on these two fluctuations are the subject of \cref{p:sec:gauge}.


\paragraph{Consistency conditions on the kernels}
\label{p:sec:hs:sextic:kernel}

The quadratic term of $\Sigma_3$ in~\eqref{p:eq:hs6:S6tilde-explicit} is not given by the sextic vertex alone.
It contains the exchange part $\sprod{2}{3} \circ \sprod{2}{2}^{-1} \circ \sprod{3}{2}$ of~\eqref{p:eq:hs6:kinetic}, two quintic vertices joined by a $\Sigma_2$ line, read with three strands in and three strands out.
The kinetic operator of $\Sigma_3$ is thus non-local in the sense of effective field theory since, unlike that of $\Sigma_2$, it contains the inverse of a vertex.
A $\Sigma_2$ line collapses to a contact term if one of its ends is the source vertex $\sprod{2}{2}(\Psi^2)$, since $\sprod{2}{2}^{-1} \circ \sprod{2}{2}$ is the identity on the quotient, and, for generic vertices, in which the exchange term is not a linear combination of the elementary vertices, only then.
  The line of $\vertex{5 \cdot 5}$ joins two quintic vertices, so it is the first that does not collapse, and this is why $\vertex{4}$ and $\vertex{5}$ needed no reduction while $\vertex{6}$ does.

The interpretation is the same one as of the quintic order.
The field $\Sigma_2$ describes more than the quartic interaction once the quintic vertex is included: it induces interactions at all orders.
The part of the sextic interaction left for $\Sigma_3$ to describe is the sextic vertex minus what $\Sigma_2$ already produces.
We can make an analogy with linear algebra: $\wprod{3}{3}$ is the quadratic form left on $\Hn{3}$ once $\Sigma_2$ is eliminated from the sextic form, that is, its Schur complement.\footnote{
    For a block matrix with blocks $A$, $B$, $B^T$ and $D$, with $D$ invertible, the Schur complement of $D$ is $A - B D^{-1} B^T$: the quadratic form left on the first block once the second is eliminated by its equation of motion.
    Here $A = \sprod{3}{3}$, $B = \sprod{2}{3}$ and $D = \tfrac{5}{6} \, \sprod{2}{2}$.
  }
  The completed square~\eqref{p:eq:hs6:DeltaS6} is built on this form.

Two conditions make $\Sigma_3$ an auxiliary field, in addition to the invertibility of $\sprod{2}{2}$ on its quotient, which the exchange part needs.
First, $\wprod{3}{3}$ must be invertible on the quotient~\eqref{p:eq:hs6:quotient}, so that $\Sigma_3$ has a propagator and is integrated out by a Gaussian integral.
Second, its kernel must couple to nothing.
Injectivity, which holds on the quotient by construction, already makes the solution~\eqref{p:eq:hs6:saddle-sigma3} unique.
Unlike the dressed operator of $\Sigma_2$, $\wprod{3}{3}$ carries no power of $\gs$: its two terms are of the same order.
Its invertibility is therefore a genuine condition on the vertex family, on the sextic vertex and on the exchange diagram together, which no expansion in $\gs$ reduces to the condition on $\sprod{2}{2}$.
It holds in the finite-dimensional models of \cref{p:app:finite-models}, with generic vertices.
For the string vertices, it is a condition on the spectrum of $\sprod{3}{3}$ on $\Hn{3}$ and on the exchange part, which can be checked level by level in level truncation.
The kernels of $\sprod{3}{3}$ and $\wprod{3}{3}$ on $\Hn{3}$ are meant, as for $\sprod{2}{2}$, in the square-integrable sense: their elements are square-integrable functions of the relative momenta at fixed total momentum.
We assume it as we did for $\sprod{2}{2}$ in \cref{p:sec:hs:quartic,p:sec:hs:quintic}: kernels in the square-integrable sense at fixed total momentum, sources of $(\wprod{3}{3})^{-1}$ and of its dressed forms lying in the image of $\wprod{3}{3}$, and kernels that stabilize as the level grows, in the sense that the quotient and the inverse have a limit at generic total momentum:
\begin{equation}
\label{p:eq:hs6:invertibility-3}
\big( \wprod{3}{3} \big)^{-1} \circ \wprod{3}{3}
  = \id
\quad \text{on } \Hn{3} / \ker \wprod{3}{3}.
\end{equation}

A kernel direction of $\wprod{3}{3}$ is a flat direction of $\Sigma_3$, with zero action.
At this order it decouples automatically, since $\Sigma_3$ enters only through $\wprod{3}{3}$.
At the next order $\Sigma_3$ enters the septic vertex.
The quotient then stays consistent only if the kernel directions are null vectors of every vertex in which $\Sigma_3$ appears, as the quintic vertex had to vanish on $\ker \sprod{2}{2}$ by~\eqref{p:eq:hs5:kernel-condition}.
The construction needs this assumption on the kernel of $\wprod{3}{3}$ from the septic order on, and we state it there (\cref{p:sec:hs:septic}).

We have taken the auxiliary fields on quotients, $\Sigma_2$ on $\Hn{2} / \ker \sprod{2}{2}$ and $\Sigma_3$ on $\Hn{3} / \ker \wprod{3}{3}$, and the action must be a function on them.
  For $\Sigma_3$ nothing is required at this order, since the field enters only through $\wprod{3}{3}$, whose kernel is quotiented.
  For $\Sigma_2$, the quintic coupling descends to the quotient if and only if the quintic vertex vanishes on the kernel,~\eqref{p:eq:hs5:kernel-condition}.
  The sextic term adds one coupling, $\dvertex{6}(\Psi \otimes \Sigma_2, \Sigma_3)$.
  Under $\Sigma_2 \to \Sigma_2 + N_2$ with $N_2 \in \ker \sprod{2}{2}$, its exchange part shifts by a term containing $\sprod{3}{2}(\Psi, N_2)$, which vanishes by~\eqref{p:eq:hs5:kernel-condition}, and its contact part by $\vertex{6}(\Psi, N_2, \Sigma_3)$.
  The cubic action therefore descends to the quotient if and only if the sextic vertex vanishes on the kernel as well:
  \begin{equation}
  \label{p:eq:hs6:kernel-condition}
  \forall N_2 \in \ker \sprod{2}{2}, \;
  \forall A \in \Hn{3}, \;
  \forall B \in \Hs:
  \quad
  \vertex{6}(N_2, A, B)
    = 0.
\end{equation}
  The quotient is thus consistent through this order when the quartic form has no kernel or, if it has one, when the quintic and sextic vertices vanish on it.
  When the vertices do not vanish on the kernel of $\sprod{2}{2}$, the quotient is not available: the quintic coupling depends on the representative of $\Sigma_2$, a complement of the kernel has to be chosen, and the eliminated action then depends on the complement, already at order $\gs^3$.
  The constraint can then be enforced by a Lagrange multiplier instead: on the kernel component only, the hybrid below, or on the whole field, the multiplier form of \cref{p:sec:all-orders}.
  The condition holds trivially at generic momentum if the quartic form has no kernel there.
  We do not try to settle here whether it holds for the actual vertex family; in principle, this can be checked level by level in level truncation.
  The quotient needs no inner product and no choice: it is canonical, and the kernel directions, which have zero action and couple to nothing, are not fields.

